\documentclass[12pt,a4paper]{article}
\usepackage{amsmath,amssymb,amsthm}
\usepackage{geometry}
\usepackage{hyperref}
\usepackage{url}
\usepackage{booktabs}
\usepackage{xcolor}
\usepackage{enumitem}
\usepackage{microtype}
\geometry{margin=2.5cm}
\hypersetup{colorlinks=true,linkcolor=blue,citecolor=blue,urlcolor=blue}
\newtheorem{theorem}{Theorem}
\newtheorem{proposition}[theorem]{Proposition}
\newtheorem{lemma}[theorem]{Lemma}
\newtheorem{conjecture}[theorem]{Conjecture}
\theoremstyle{definition}
\newtheorem{definition}[theorem]{Definition}
\newtheorem{remark}[theorem]{Remark}
\newcommand{\ex}{\mathrm{ex}}

\title{\textbf{$C_4$-Free Subgraphs of the Hypercubes $Q_6$, $Q_7$, and $Q_8$:\\
Odd Squares, Fully Frustrated Models, and Computational Structure}}
\author{Minamo Minamoto\\[4pt]
\small Independent\\
\small \texttt{minamominamoto4f5683f6@gmail.com}}
\date{March 2026 (revised August 2026)}

\begin{document}
\maketitle

\begin{center}
\small This is a revision of a manuscript first circulated as
arXiv:2603.29127 (v1--v4, March--May 2026), under the title
``New Lower Bounds for $C_4$-Free Subgraphs of the Hypercubes $Q_6$,
$Q_7$, and $Q_8$: Constructions, Structure, and Computational
Method.'' That earlier version announced $\ex(Q_8,C_4)\ge680$ and did
not include the odd-square material of Section~\ref{sec:oddsquare};
see the Acknowledgements for the circumstances of the revision.
\end{center}
\medskip

\begin{abstract}
We study $C_4$-free subgraphs of the hypercubes $Q_6$, $Q_7$, and $Q_8$,
giving explicit lower bounds $\ex(Q_7,C_4)\ge304$ and $\ex(Q_8,C_4)\ge682$,
together with a computational reproduction of the 132-edge lower-bound
construction for $\ex(Q_6,C_4)=132$, combined with the known upper bound
of Harborth and Nienborg~\cite{HN94} (we did not independently verify
ILP optimality for the upper bound within a practical runtime; see
Section~\ref{sec:q6}).
After this manuscript was first circulated, a reader (S.~Lai) identified an
overlooked connection to a line of statistical-physics literature on fully
frustrated hypercubes: under the standard correspondence between edge sets
meeting every $4$-cycle (\emph{square}) an odd number of times and fully
frustrated signed hypercubes, the 1979 construction of Derrida, Pomeau,
Toulouse, and Vannimenus~\cite{DPTV79} already yields a 304-edge $C_4$-free
subgraph of $Q_7$, and the $D{=}8$ ground state reported by Marinari, Parisi,
and Ritort~\cite{MPR95} yields a 682-edge $C_4$-free subgraph of $Q_8$ that
improves our originally announced 680-edge construction. For $Q_7$, this
value coincides with one of our own previously found 304-edge solutions,
which we confirm independently satisfies the odd-square condition. For
$Q_8$, we derived and machine-verified an explicit 682-edge witness
attaining the bound MPR report attaining, using their canonical fully
frustrated coupling (derivation documented in
Section~\ref{sec:oddsquareverify}); this witness is not claimed to be a
reconstruction of any particular configuration MPR may have used
internally, since they did not publish one. This 682-edge
witness is presented here as the paper's headline lower bound in place
of the earlier 680-edge witness.

Writing $g(n)$ for the maximum size of an \emph{odd-square} edge set (every
square meets it in exactly one or three edges), we show $g(6)=132$,
$g(7)=304$, and $g(8)=682$ exactly. For $n=7,8$ the upper bound follows
from the local-field argument of Derrida et al.~\cite{DPTV79} and the
lower bound from explicit constructions; for $n=6$ we instead combine an
explicit odd-square construction already given in~\cite{DPTV79} itself
(lower bound $g(6)\ge132$) with the independent exact value
$\ex(Q_6,C_4)=132$ of Harborth and Nienborg~\cite{HN94}, which bounds
$g(6)$ from above since every odd-square set is $C_4$-free. This settles
the odd-square subclass completely for $n=6,7,8$ without resolving the
unrestricted problem $\ex(Q_n,C_4)$, since a general $C_4$-free subgraph
need not be odd-square.

An exhaustive audit of the $19\,866$ previously published $304$-edge $Q_7$
solutions shows that only $389$ of them lie in the odd-square subclass; the
remaining $19\,477$ do not. The odd-square construction therefore does
not subsume or fully explain the structural classification given here
(one profile type, Type~18 with 101 solutions, is entirely odd-square,
but the other 19 types are not): the $19\,866$ solutions'
common structural core (degree sequence $\{4^{32},5^{96}\}$, spectral radius
in $[4.78543,4.79129]$, local maximality) and their classification into 20
dimension-profile types are independent results. Whether these $19\,866$
solutions exhaust all 304-edge $C_4$-free subgraphs of $Q_7$ remains open.

For $Q_8$, the reconstructed 682-edge odd-square witness is locally maximal
in a strong sense: all 342 non-edges of $Q_8$ create at least 3 new $C_4$s
upon addition (compared with the earlier 680-edge simulated-annealing
Solution~B, for which 2 of 344 non-edges created only 1 or 2 new $C_4$s;
see Section~\ref{sec:q8} for both 680-edge solutions released). All
bounds are witnessed by explicit edge-list constructions with
machine-verifiable certificates: a complete, deterministic enumeration of
all four-cycles (240 for $Q_6$, 672 for $Q_7$, 1\,792 for $Q_8$) establishes
$C_4$-freeness, and every certificate reported here, including the
odd-square reconstructions, is reproducible by an independent,
dependency-free verifier.

The original ($Q_7$, $Q_8$) constructions are found by a two-phase simulated
annealing algorithm, described in full detail; the released search code
computes an $\mathrm{Aut}(Q_n)$-diversified restart point before each
trial but does not pass it into the phase-1 search, which always
restarts from a fresh uniform-random edge set instead (see
Section~\ref{sec:method}). For $Q_6$ we additionally give an ILP formulation whose
feasible value matches the known exact value $\ex(Q_6,C_4)=132$; we did
not independently verify the ILP's optimality within a practical
runtime, and the upper bound itself rests on the combinatorial proof
of Harborth and Nienborg~\cite{HN94}, which we reproduce on the lower
(construction) side.

Edge lists, ILP files, the odd-square reconstruction and audit scripts, and
source code are made available at\\
\url{https://github.com/minamominamoto/c4free-hypercube}.
%% ZERO: as of this revision, the repository has not yet been updated to
%% include the odd-square reconstruction/audit scripts, q8_odd_square_682.json,
%% or this revision's edge lists; its README still states ex(Q8,C4)>=680.
%% The updated files (README.md, this tex/pdf, all scripts and data) have
%% been prepared as a flat set ready to add to your local clone and push
%% yourself -- git push credentials are not available in this sandbox.
%% Ideally push as a tagged/archived release with a SHA-256 manifest,
%% before submission, so this claim is true at the time a reader follows
%% the link. Remove this comment once done.
\end{abstract}

\medskip
\noindent\textbf{MSC 2020:} 05C35, 05C62, 05C50.\\
\textbf{Keywords:} hypercube, $C_4$-free, extremal graph theory,
simulated annealing, integer linear programming, dimension profile.

\tableofcontents

%% ============================================================
\section{Introduction}
%% ============================================================

The $n$-dimensional hypercube $Q_n$ has vertex set $\{0,1\}^n$, two vertices
adjacent iff they differ in exactly one coordinate; it has $2^n$ vertices and
$n\cdot2^{n-1}$ edges.  The problem of determining
$\ex(Q_n,C_4)=\max\{|E(G)|:G\subseteq Q_n,\,G\text{ is }C_4\text{-free}\}$
was raised by Erd\H{o}s~\cite{Erd84} and remains open in general.

The best known asymptotic bounds are
\begin{equation}\label{eq:bounds}
\tfrac{1}{2}(n+\sqrt{n})\cdot2^{n-1}\;\le\;\ex(Q_n,C_4)\;\le\;(0.60318+o(1))\cdot n\cdot2^{n-1},
\end{equation}
where the lower bound (valid for $n=4^r$) is due to Brass, Harborth, and
Nienborg~\cite{BHN95}. The upper bound is a statement about the edge
Tur\'an \emph{density} $\pi_e=\limsup_n\ex(Q_n,C_4)/(n2^{n-1})$, not a
pointwise inequality valid at every finite $n$; the constant $0.60318$
is due to Baber~\cite{Bab12}, who used the semidefinite flag algebra
method to improve the earlier bound of Thomason and Wagner~\cite{TW09};
the intermediate value $0.6068$ was obtained independently by
Baber~\cite{Bab12} and by Balogh, Hu, Lidick\'y, and Liu~\cite{BHLL12}.
For general $n\ge9$, the weaker bound
$\tfrac{1}{2}(n+0.9\sqrt{n})\cdot2^{n-1}$ applies~\cite{BHN95}.
Exact values have been determined for small $n$; Harborth and Nienborg~\cite{HN94}
proved $\ex(Q_6,C_4)=132$.

Our main results are:

\begin{theorem}\label{thm:main}
\begin{enumerate}[label=(\roman*)]
  \item $\ex(Q_7,C_4)\ge304$.
  \item $\ex(Q_8,C_4)\ge682$.
\end{enumerate}
\end{theorem}

Part (i) is not new: it follows from the 1979 construction of Derrida,
Pomeau, Toulouse, and Vannimenus~\cite{DPTV79} (Section~\ref{sec:oddsquare}),
and we present it here as an independently reproduced and machine-verified
result rather than as a novel bound. Part (ii) improves our originally
announced bound $\ex(Q_8,C_4)\ge680$, but here too the bare existence
claim is not new in substance: Marinari, Parisi, and Ritort~\cite{MPR95}
already report having exhibited a $D=8$ ground state attaining Derrida
et al.'s improved energy lower bound, and combined with the
signed-graph/odd-square correspondence and edge-count--field-sum
translation we give in Section~\ref{sec:oddsquare}, their published
existence claim already implies the existence of a 682-edge odd-square
subgraph of $Q_8$, independently of any witness search of our own. What
Part~(ii) contributes beyond that prior existence claim is: an explicit
graph-theoretic translation of MPR's result into edge-count terms
(Section~\ref{sec:oddsquare-exact}); a publicly released, explicit,
machine-verifiable 682-edge certificate, which MPR's paper does not
supply (they report the result in energy units only, without an
explicit spin configuration); and the structural analysis of that
certificate (Proposition~\ref{prop:q8oddsquare}). The specific
certificate we release attains the same bound value MPR report, but its
own provenance is not fully documented (see the note in
Section~\ref{sec:oddsquareverify}); we do not claim it reproduces any
particular configuration MPR may have used internally.
Both bounds are realised by explicit constructions (supplementary files
\texttt{q7\_edges\_304.jsonl.part1--3} and \texttt{q8\_odd\_square\_682.json};
the originally announced 680-edge $Q_8$ witness, \texttt{q8\_edges\_680.jsonl},
is retained as a simulated-annealing data point, see
Section~\ref{sec:q8}).

Within the \emph{odd-square} subclass --- edge sets meeting every square in
an odd number of edges --- the extremal value $g(n)$ is determined exactly
for $n=7,8$:

\begin{theorem}\label{thm:oddsquare}
$g(6)=132$, $g(7)=304$, and $g(8)=682$.
\end{theorem}

\noindent Theorem~\ref{thm:oddsquare} is proved in
Section~\ref{sec:oddsquare}. It does not resolve
Theorem~\ref{thm:main}'s unrestricted bounds, since $g(n)\le\ex(Q_n,C_4)$
and a general $C_4$-free subgraph need not be odd-square: of the $19\,866$
published 304-edge $Q_7$ solutions, only 389 are odd-square
(Section~\ref{sec:structure}).

\begin{table}[ht]
\centering
\caption{Known values and lower bounds of $\ex(Q_n,C_4)$.
  The BHN column gives Brass--Harborth--Nienborg estimates from
  Eq.~\eqref{eq:bounds}; that equation's tight form applies only for
  $n=4^r$ and its weaker form only for $n\ge9$~\cite{BHN95}, so the
  $n=5,6$ entries are not covered by either stated form (marked n/a)
  and the $n=7,8$ entries are extrapolations of the $n\ge9$ form
  \emph{beyond} its stated domain, included for illustrative comparison
  only, not as a proven bound.}
\label{tab:values}\medskip
\begin{tabular}{ccccl}
\toprule
$n$ & $|E(Q_n)|$ & $\ex(Q_n,C_4)$ & BHN (Eq.~\ref{eq:bounds}) & Source\\
\midrule
4 &   32 & 24  & $24$ (tight form, $n=4^1$) & Chung~\cite{Chu92}\\
5 &   80 & 56  & n/a (outside stated domain) & Emamy-K et al.~\cite{EMR92}\\
6 &  192 & 132 & n/a (outside stated domain) & Harborth--Nienborg~\cite{HN94} (reproduced)\\
7 &  448 & $\ge\mathbf{304}$ & $\approx300.2$ (extrapolated) & Derrida et al.~\cite{DPTV79} (reproduced)\\
8 & 1024 & $\ge\mathbf{682}$ & $\approx674.9$ (extrapolated) & witness attaining bound of~\cite{MPR95} (this paper)\\
\bottomrule
\end{tabular}
\end{table}

The paper is organised as follows.
Section~\ref{sec:verify} describes the $C_4$ enumeration and verification
procedure used throughout.
Sections~\ref{sec:q7}--\ref{sec:q8} present the originally announced $Q_7$
and $Q_8$ simulated-annealing constructions.
Section~\ref{sec:oddsquare} introduces the odd-square subclass, proves
Theorem~\ref{thm:oddsquare}, and presents the reconstructed 682-edge $Q_8$
witness and its verification.
Section~\ref{sec:structure} classifies the $19\,866$ $Q_7$ solutions found
by their dimension profiles, yielding exactly 20 types, and reports the
odd-square audit (389 of 19\,866).
Section~\ref{sec:method} describes the two-phase simulated annealing algorithm.
Section~\ref{sec:q6} reproduces $\ex(Q_6,C_4)=132$ computationally on the
lower-bound side, with an ILP formulation for the upper-bound side.
Section~\ref{sec:open} lists open problems.

%% ============================================================
\section{$C_4$ Enumeration and Verification}\label{sec:verify}
%% ============================================================

\begin{lemma}[$C_4$ enumeration in $Q_n$]\label{lem:c4}
Every four-cycle of $Q_n$ is uniquely determined by a pair of coordinate
directions $0\le i<j\le n-1$ and a base vertex $b\in\{0,1\}^n$ with
$b_i=b_j=0$.  Its four vertices are
\[
  b,\quad b\oplus e_i,\quad b\oplus e_j,\quad b\oplus e_i\oplus e_j,
\]
and the total number of four-cycles in $Q_n$ is $\binom{n}{2}\cdot2^{n-2}$.
For $n=6$: $240$; for $n=7$: $672$; for $n=8$: $1\,792$.
\end{lemma}

\noindent\textbf{Verification algorithm.}\quad
Given a candidate edge set $E\subseteq E(Q_n)$ stored as a hash set,
the following procedure performs a complete, deterministic $C_4$-freeness check:
\begin{quote}\small
\textbf{for each} pair $(i,j)$ with $0\le i<j\le n-1$:\\
\quad\textbf{for each} $b\in\{0,1\}^n$ with $b_i=b_j=0$:\\
\quad\quad let $v_1=b,\;v_2=b\oplus e_i,\;v_3=b\oplus e_j,\;v_4=b\oplus e_i\oplus e_j$\\
\quad\quad \textbf{if} $\{v_1v_2,\,v_1v_3,\,v_2v_4,\,v_3v_4\}\subseteq E$:
  \textbf{return} \textsc{violation}\\
\textbf{return} \textsc{c4-free}
\end{quote}
We applied this algorithm to all constructions and all $19\,866$ solutions
of $Q_7$; every call returned \textsc{c4-free}.

%% ============================================================
\section{The $Q_7$ Construction}\label{sec:q7}
%% ============================================================

We constructed a $C_4$-free subgraph of $Q_7$ on 304 edges by
penalty-based simulated annealing (Section~\ref{sec:method}),
then certified it by Lemma~\ref{lem:c4}'s algorithm.

\begin{proposition}\label{prop:q7}
The 304-edge construction has the following properties:
\begin{enumerate}[label=(\roman*)]
  \item Degree sequence: 32 vertices of degree 4 and 96 of degree 5
    (degree sum $608=2\cdot304$).
  \item Spectral radius $\lambda_1\approx4.787$ (a genuine computed
    quantity for this construction); $\lambda_n\approx-4.787$, i.e.\
    $\lambda_1+\lambda_n=0$ to machine precision, which is automatic
    for any subgraph of the bipartite $Q_7$ and serves here only as a
    sanity check on the eigenvalue computation, not as an independent
    structural finding.
  \item $\mathrm{Tr}(A^4)=5216$, attaining the $C_4$-free trace equality
    $\sum_i\deg(i)^2+\sum_{uv\in E}(\deg u+\deg v)-2|E|$.
  \item Local maximality: every non-edge of $Q_7$ creates a $C_4$ when added.
\end{enumerate}
\end{proposition}

Through $19\,866$ independent runs (Section~\ref{sec:landscape})
we found $19\,866$ $C_4$-free subgraphs of $Q_7$ on 304 edges with
pairwise distinct edge sets (not quotiented by $\mathrm{Aut}(Q_7)$),
all sharing the structural properties stated in Proposition~\ref{prop:rigid}.

%% ============================================================
\section{The $Q_8$ Construction}\label{sec:q8}
%% ============================================================

This section documents the originally announced $Q_8$ witness, obtained
independently of the odd-square subclass by the same simulated-annealing
method as the $Q_7$ construction of Section~\ref{sec:q7}. It has since been
superseded as the paper's headline lower bound by the 682-edge odd-square
reconstruction of Section~\ref{sec:oddsquare}; it is retained here as an
independent structural data point and as a demonstration of the
simulated-annealing method on its own terms.

Simulated annealing produced two $C_4$-free subgraphs of $Q_8$ on 680
edges (both released in \texttt{q8\_edges\_680.jsonl}), each certified by
exhaustive enumeration of all $1\,792$ four-cycles. We refer to them as
Solution~A (first line of the file) and Solution~B (second line); the
two are structurally distinct and are kept separate below, since an
earlier draft of this section inadvertently described properties of
Solution~A and Solution~B as if they belonged to a single construction.

\begin{proposition}\label{prop:q8struct}
Both 680-edge solutions are $C_4$-free and locally maximal (every
non-edge of $Q_8$ creates at least one $C_4$ upon addition), with edge
density $680/1024=66.4\%$ (compared with $304/448=67.9\%$ for the $Q_7$
construction) and $\lambda_1+\lambda_n=0$ (automatic for any subgraph of
the bipartite $Q_8$; not by itself a distinguishing structural feature).
They differ as follows:
\begin{enumerate}[label=(\roman*)]
  \item Solution~A: degree sequence $\{4^8,\,5^{160},\,6^{88}\}$
    (degree sum $1360=2\cdot680$); square-intersection histogram
    $\{1\!:\!308,\,3\!:\!1484\}$ over the $1\,792$ squares (incidentally
    odd-square, though not constructed as such); non-edge violation
    distribution $\{2\!:\!3,\,3\!:\!48,\,4\!:\!144,\,5\!:\!136,\,6\!:\!13\}$
    (minimum $2$ new $C_4$s per non-edge).
  \item Solution~B: degree sequence $\{4^7,\,5^{162},\,6^{87}\}$
    (degree sum $1360=2\cdot680$); square-intersection histogram
    $\{1\!:\!302,\,2\!:\!12,\,3\!:\!1478\}$ (not odd-square: $12$ squares
    meet it in exactly $2$ edges); non-edge violation distribution
    $\{1\!:\!1,\,2\!:\!1,\,3\!:\!49,\,4\!:\!153,\,5\!:\!124,\,6\!:\!16\}$
    (minimum $1$ new $C_4$, attained by exactly one non-edge).
\end{enumerate}
\end{proposition}

\begin{remark}
Brass--Harborth--Nienborg's tight bound $\tfrac12(n+\sqrt n)\cdot2^{n-1}$
applies only for $n=4^r$, and their weaker bound
$\tfrac12(n+0.9\sqrt n)\cdot2^{n-1}$ is stated only for $n\ge9$
(Eq.~\eqref{eq:bounds},~\cite{BHN95}); neither covers $n=8$. Evaluating
the $n\ge9$ form at $n=8$ nonetheless gives $\approx674.9$; both
680-edge constructions and the 682-edge odd-square reconstruction of
Section~\ref{sec:oddsquare} exceed this extrapolated figure (by 5.1 and
7.1 edges respectively), but we do not claim this as a comparison
against a proven bound.
\end{remark}

\subsection{Local optimality of the 680-edge constructions}

Solution~B is locally maximal in the weakest possible sense among the
two: among its $344$ non-edges, exactly $1$ creates exactly $1$ new
$C_4$; exactly $1$ creates exactly $2$; the remaining $342$ create $3$
or more (Proposition~\ref{prop:q8struct}(ii)). Solution~A's non-edges
all create at least $2$ new $C_4$s.

Over $1\,076$ independent simulated-annealing trials targeting $681$
edges, the minimum $C_4$ violation count observed was $1$ (never $0$),
where the violation count denotes the number of four-cycles fully
contained in the candidate edge set. This was originally reported as
computational evidence for the conjecture $\ex(Q_8,C_4)=680$; that
conjecture is withdrawn, since the odd-square reconstruction of
Section~\ref{sec:oddsquare} gives $\ex(Q_8,C_4)\ge682>680$. Since a
stochastic annealing run does not define a graph-theoretic neighbourhood
search (no fixed radius or reachability guarantee is established), the
681-edge search failure is evidence only about the behaviour of the
stated heuristic across $1\,076$ independent fresh-random-restart trials
targeting 681 edges (Section~\ref{sec:method}), not a mathematical
statement about $\ex(Q_8,C_4)$ or about any formally defined
neighbourhood of the 680-edge solutions.

Table~\ref{tab:compare} compares the $Q_7$ construction with Solution~B
(the 680-edge solution whose local-optimality numbers are discussed
above); see Table~\ref{tab:oddsquare} in Section~\ref{sec:oddsquare} for
the odd-square reconstructions.

\begin{table}[ht]
\centering
\caption{Structural comparison of the $Q_7$ construction and $Q_8$
  Solution~B (see Proposition~\ref{prop:q8struct} for Solution~A).}
\label{tab:compare}\medskip
\begin{tabular}{lcc}
\toprule
Parameter & $Q_7$ (304 edges) & $Q_8$ Solution~B (680 edges)\\
\midrule
Total edges $|E(Q_n)|$ & 448 & 1024\\
Achieved edges & 304 & 680\\
Edge density & 67.9\% & 66.4\%\\
Min degree & 4 & 4\\
Max degree & 5 & 6\\
Degree sequence & $\{4^{32},5^{96}\}$ & $\{4^7,5^{162},6^{87}\}$\\
BHN (Eq.~\ref{eq:bounds}), extrapolated & $\approx300.2$ & $\approx674.9$\\
Improvement & $+3.8$ & $+5.1$\\
\bottomrule
\end{tabular}
\end{table}

%% ============================================================
\section{The Odd-Square Subclass}\label{sec:oddsquare}
%% ============================================================

\subsection{Definition and equivalence to fully frustrated hypercubes}

\begin{definition}
An edge set $E\subseteq E(Q_n)$ is \emph{odd-square} if every square
(four-cycle) of $Q_n$ meets $E$ in an odd number of edges, i.e.\ one or
three of its four edges lie in $E$. Write
$g(n)=\max\{|E|:E\text{ is odd-square in }Q_n\}$.
\end{definition}

Every odd-square edge set is $C_4$-free, since no square can meet it in all
four edges; hence $g(n)\le\ex(Q_n,C_4)$. The odd-square condition is exactly
the condition studied under a different name in the statistical-physics
literature on \emph{fully frustrated hypercubes}: assign to each edge $e$ a
sign $\sigma(e)=+1$ if $e\in E$ and $\sigma(e)=-1$ otherwise; the square
condition becomes ``the product of signs around every square is $-1$'', i.e.\
every square is frustrated.

\begin{lemma}\label{lem:switching}
Let $\sigma_1,\sigma_2:E(Q_n)\to\{\pm1\}$ both have product $-1$ on every
square. Then there is a function $t:\{0,1\}^n\to\{\pm1\}$ with
$\sigma_2(x,y)=\sigma_1(x,y)\,t(x)t(y)$ for every edge $(x,y)$.
\end{lemma}
\begin{proof}
Let $\tau=\sigma_1\sigma_2$ (pointwise product). On every square $C$,
$\prod_{e\in C}\tau(e)=\bigl(\prod_{e\in C}\sigma_1(e)\bigr)\bigl(\prod_{e\in C}\sigma_2(e)\bigr)=(-1)(-1)=+1$.
The map sending a cycle $C$ to $\prod_{e\in C}\tau(e)\in\{\pm1\}$ is a group
homomorphism from the cycle space of $Q_n$ (over $\mathrm{GF}(2)$, under
symmetric difference) to $\{\pm1\}$, since edges shared by two cycles cancel
in their symmetric difference and likewise cancel in the product. The
squares of $Q_n$ generate its cycle space (a standard fact for hypercubes:
every cycle is a symmetric difference of squares), so this homomorphism is
determined by its values on the squares; since it is $+1$ on every square,
it is $+1$ on every cycle. A signature with product $+1$ on every cycle of
a connected graph is a \emph{coboundary}: $\tau(x,y)=t(x)t(y)$ for some
$t:V\to\{\pm1\}$ (the classical balance theorem for signed graphs, e.g.\
via Harary's theorem). Solving for $\sigma_2$ gives the claim.
\end{proof}

Lemma~\ref{lem:switching} says any two odd-square signatures are related by
a \emph{vertex switching} $t$; in Ising terms, replacing $s(x)$ by $t(x)s(x)$
for all $x$ is a gauge transformation that preserves the ground-state
energy (equivalently, $|E_2|=|\{e:\sigma_2(e)=+1\}|$ ranges over exactly
the same set of values as $|E_1|$ does, as $\sigma_1$ ranges over a full
switching class, when $\sigma_2$ ranges over all odd-square signatures).
Consequently, $g(n)$ — defined as a maximum over \emph{all} odd-square
edge sets — equals the maximum of $|E|$ over the switching class of any
one fixed fully-frustrated signature, in particular the reference coupling
$J$ used in Section~\ref{sec:oddsquareverify}'s witness; optimising that
one fixed $J$ (or, for $D=8$, exhibiting a ground state attaining
Derrida et al.'s improved energy lower bound for it, as
Marinari--Parisi--Ritort report~\cite{MPR95}) therefore does optimise
over the class defining $g(n)$, not merely over a possibly-proper
subclass of it.

Under the further correspondence between signed
hypercubes and Ising spin configurations $s:\{0,1\}^n\to\{\pm1\}$ via
$\sigma(x,y)=J(x,y)\,s(x)s(y)$ for a fixed reference coupling $J$ with every
square frustrated, maximising $|E|$ (i.e.\ maximising the number of
\emph{positive} edges, equivalently minimising the number of negative
edges) is exactly the problem of minimising the ground-state energy of the
fully frustrated Ising hypercube studied by Derrida, Pomeau, Toulouse, and
Vannimenus~\cite{DPTV79} and, for $D=8$, by Marinari, Parisi, and
Ritort~\cite{MPR95}: the ground-state energy is (up to an affine rescaling)
the number of negative edges, $|E(Q_n)|-|E|$. In the language of signed
graphs, $|E(Q_n)|-g(n)$ is the \emph{frustration index} of the fully
frustrated signature on $Q_n$; see also~\cite{LNSW25} for a recent treatment
of this quantity in the context of Huang's proof of the sensitivity
conjecture. This sign convention matches the one used by the reconstruction
script \texttt{generate\_q8\_682.py} (Section~\ref{sec:oddsquareverify}),
where $E$ is exactly the set of edges with $J(x,\mathrm{dim})\,s(x)s(y)=+1$;
for a vertex $v$ of degree $\deg_E(v)$ in $E$, we call
$h(v)=2\deg_E(v)-n$ its \emph{local field} under this convention (matching
the local-field histogram reported in Proposition~\ref{prop:q8oddsquare}
below). We do not reproduce Derrida et al.'s derivation of the local-field
lower bound here; Section~\ref{sec:oddsquare-exact} cites it as an external
result, established for $D\le7$ by explicit construction and used by
Marinari--Parisi--Ritort's $D=8$ construction, which we reconstruct and
verify independently in Section~\ref{sec:oddsquareverify} without relying on
the bound's derivation.

This connection was brought to our attention, after the original circulation
of this manuscript, by S.~Lai (see Acknowledgements), who identified the
relevance of~\cite{DPTV79,MPR95} and located the improved $Q_8$ result.

\subsection{Exact values for \texorpdfstring{$n=6,7,8$}{n=6,7,8}}\label{sec:oddsquare-exact}

Derrida et al.~\cite{DPTV79} give a local-field lower bound on the number of
negative (frustrated, in their terminology) edges, tight enough to
determine the fully frustrated hypercube's ground-state energy exactly
for $D\le7$ via an explicit recursive construction, and conjectured it
remains tight for $D\ge8$. For $D=6$, they give this explicit
construction directly: ``this configuration for $H_6$ contains 60
negative bonds (out of a total of 192)''~\cite[p.~622]{DPTV79}, i.e.\
$192-60=132$ positive (unfrustrated) edges, with the construction
guaranteed fully frustrated (every square meeting it in an odd number
of negative edges) by their general recursive method. Under the
correspondence of Section~\ref{sec:oddsquare} this is an explicit
odd-square witness with $132$ edges, so $g(6)\ge132$; combined with
$g(6)\le\ex(Q_6,C_4)=132$~\cite{HN94}, this gives $g(6)=132$. We did not
reconstruct this $D=6$ configuration ourselves (unlike the $D=7,8$
cases below); the value $132$ rests on the quoted primary-source count
together with the independently proved exact value from~\cite{HN94}, not
on our own machine verification of the specific bond configuration.

Marinari, Parisi, and Ritort~\cite{MPR95} report having exhibited, by
simulated annealing, an explicit $D=8$ ground state attaining Derrida et
al.'s improved energy lower bound; their paper reports the result in
energy units and does not include the explicit spin configuration or its
translation into edge-count terms. We make this translation explicit and
independently checkable. Writing $U$ for one part of the bipartition of
$Q_n$ ($|U|=2^{n-1}$) and $h(v)=2\deg_E(v)-n$ for the local field of
Section~\ref{sec:oddsquare} restricted to $E$, a bipartite double-counting
identity gives $\sum_{v\in U}h(v)=2|E|-n\cdot2^{n-1}$ for \emph{any} edge
set $E$ (odd-square or not), hence
\[
|E|=\frac{n\cdot2^{n-1}+\sum_{v\in U}h(v)}{2}.
\]
For $n=7$, Derrida et al.\ both derive the field-sum bound and prove it
is achieved by an explicit construction (their $d\le7$ case): the bound
value is $\sum_{v\in U}h(v)=160$ (table~I of~\cite{DPTV79}), giving
$|E|=(448+160)/2=304$. For $n=8$ the situation is different: Derrida et
al.\ give the same bound value, $\sum_{v\in U}h(v)=340$, giving
$|E|=(1024+340)/2=682$, but do \emph{not} themselves establish that this
bound is achievable — they state only a conjecture for $D\ge8$, without
an explicit construction. Achievability at $D=8$ is exactly what
Marinari, Parisi, and Ritort~\cite{MPR95} later established, by
exhibiting a ground state attaining the bound. We stress what this checks and what it does not:
since the identity above holds for \emph{any} edge set, computing
$\sum_{v\in U}h(v)$ from a witness already known to have $304$ or $682$
edges automatically returns $160$ or $340$ respectively — we confirmed
this (for a representative odd-square 304-edge $Q_7$ solution,
Section~\ref{sec:oddsquareaudit}, and for the 682-edge $Q_8$ witness
below), but it is a consistency check on the witnesses, not an
independent verification that $160$/$340$ are the \emph{maximal}
achievable field sums. That upper-bound claim rests entirely on Derrida
et al.'s local-field argument and Table~I~\cite{DPTV79}, which we cite
as an external result (Section~\ref{sec:oddsquare-exact}) rather than
re-derive here. Under the correspondence of
Section~\ref{sec:oddsquare}, attaining Derrida et al.'s bound at $D=8$
therefore corresponds to $g(8)=682$; the explicit 682-edge witness
realising it is due to the present reconstruction below, prompted by
S.~Lai's identification of the connection to~\cite{DPTV79,MPR95}.
Together with the $D\le7$ case this gives Theorem~\ref{thm:oddsquare}:
$g(6)=132$, $g(7)=304$, and $g(8)=682$.

\subsection{Verification}\label{sec:oddsquareverify}

We reconstructed explicit witnesses for both values and machine-verified
them independently of~\cite{DPTV79,MPR95}.

\textbf{$Q_7$, $g(7)=304$.} The 304-edge odd-square value coincides with our
originally announced $Q_7$ lower bound (Theorem~\ref{thm:main}(i)); of the
$19\,866$ solutions we found by simulated annealing (Section~\ref{sec:q7}),
389 are themselves odd-square witnesses of $g(7)=304$
(Section~\ref{sec:structure}), independently confirming reachability of this
value without reference to~\cite{DPTV79}.

\textbf{$Q_8$, $g(8)=682$.} We obtained a 682-edge odd-square witness
for $Q_8$ from an explicit $256$-bit spin configuration together with the
canonical fully frustrated coupling
$J(x,\mathrm{dim})=(-1)^{\mathrm{popcount}(x\,\&\,(2^{\mathrm{dim}}-1))}$
(one coupling convention realising a fully frustrated signature; see
Section~\ref{sec:oddsquare} supplementary script \texttt{generate\_q8\_682.py}).
The script is dependency-free (Python standard library only),
recomputes the edge set from the spin configuration, and asserts every
structural quantity below against the recomputed data before writing the
certificate; it does not take any of these quantities as input.
\emph{Provenance of the spin configuration.} The 256-bit configuration
was derived by the author on 17--18 August 2026, using the canonical
fully frustrated coupling of Marinari, Parisi, and Ritort~\cite{MPR95}
(private correspondence with S.~Lai, who independently reconstructed the
identical 682-edge set from the shared spin vector and cross-checked all
1,792 squares, the degree distribution, and the local-field distribution
before this material was incorporated into the manuscript). We do not
claim this configuration is identical to any specific configuration MPR
may have used internally, since they did not publish one; it is an
independently derived witness attaining the same bound.

\begin{proposition}\label{prop:q8oddsquare}
The witness edge set satisfies:
\begin{enumerate}[label=(\roman*)]
  \item $682$ edges, spanning all $256$ vertices of $Q_8$.
  \item Square-intersection histogram $\{1\!:\!301,\;3\!:\!1491\}$ over
    all $1\,792$ squares (odd-square condition holds everywhere; zero
    squares meet the edge set in $0$, $2$, or $4$ edges).
  \item Degree sequence $\{4^2,\,5^{168},\,6^{86}\}$; degree sum
    $1364=2\cdot682$.
  \item Local-field histogram $\{0\!:\!2,\,2\!:\!168,\,4\!:\!86\}$, where
    the local field at a vertex is (positive degree) $-$ (negative degree)
    under the signed-graph correspondence above.
  \item Local maximality: every one of the $342$ non-edges of $Q_8$
    creates at least $3$ new $C_4$s upon addition (distribution
    $\{3\!:\!41,\,4\!:\!159,\,5\!:\!120,\,6\!:\!22\}$), a strictly stronger
    local-maximality margin than the originally announced 680-edge
    Solution~B of Section~\ref{sec:q8}, for which 2 of 344 non-edges
    created only 1 or 2 new $C_4$s.
\end{enumerate}
\end{proposition}

All quantities in Proposition~\ref{prop:q8oddsquare} were recomputed
independently of the reconstruction script by direct enumeration of all
$1\,792$ squares against the released edge list, using the strict
definition of $C_4$-freeness from Lemma~\ref{lem:c4} (no square may
have all four edges present); this independent check found zero violations.
The reconstruction, its SHA-256 checksum, and an independent audit were
cross-checked against S.~Lai's separately produced implementation, with
full agreement.

Table~\ref{tab:oddsquare} summarises the odd-square reconstructions
alongside the corresponding simulated-annealing witnesses of
Sections~\ref{sec:q7}--\ref{sec:q8}.

\begin{table}[ht]
\centering
\caption{Odd-square reconstructions compared with the simulated-annealing
  witnesses of Sections~\ref{sec:q7}--\ref{sec:q8}.}
\label{tab:oddsquare}\medskip
\begin{tabular}{lcc}
\toprule
Parameter & $Q_7$ odd-square & $Q_8$ odd-square\\
\midrule
Edges (value of $g(n)$) & 304 & 682\\
Source & one of 389 SA solutions & witness attaining bound of~\cite{MPR95}\\
Non-odd-square SA witness (Solution~B, this paper) & --- & 680 (Section~\ref{sec:q8})\\
Improvement over Solution~B & --- & $+2$\\
\bottomrule
\end{tabular}
\end{table}

%% ============================================================
\section{Structural Classification of $Q_7$ Solutions}\label{sec:structure}
%% ============================================================

\subsection{Dimension profiles and the twenty types}\label{sec:profiles}

\begin{definition}[Dimension profile]
For $G\subseteq Q_7$, let $e_i(G)=|\{uv\in E(G):u\oplus v=2^i\}|$.
The \emph{dimension profile} is the sorted tuple
$\pi(G)=(e_{\sigma(0)}\ge\cdots\ge e_{\sigma(6)})$.
\end{definition}

The dimension profile is invariant under coordinate permutations (a subgroup
of $\mathrm{Aut}(Q_7)$ of order $7!=5040$), giving a coarse but computable
classification.

\begin{proposition}\label{prop:twenty}
Among the $19\,866$ solutions, the dimension profile takes exactly
$\mathbf{20}$ distinct values.
\end{proposition}

Table~\ref{tab:types} lists all 20 types with frequencies and structural invariants.

\begin{table}[ht]
\centering\small
\caption{The 20 dimension-profile types of 304-edge $C_4$-free subgraphs of
  $Q_7$, sorted by decreasing frequency.
  $H$: Shannon entropy of normalised profile (bits).}
\label{tab:types}\medskip
\begin{tabular}{clrrl}
\toprule
Rank & Profile $\pi$ & Solutions & $H$ & $\lambda_1$ (mean)\\
\midrule
 1 & $[44,44,44,44,43,43,42]$ & 3155 & 2.8072 & 4.787\\
 2 & $[45,45,45,43,43,42,41]$ & 2913 & 2.8065 & 4.788\\
 3 & $[45,45,45,43,43,43,40]$ & 2200 & 2.8063 & 4.787\\
 4 & $[44,44,44,44,44,43,41]$ & 2116 & 2.8069 & 4.787\\
 5 & $[46,46,46,42,42,42,40]$ & 1756 & 2.8053 & 4.787\\
 6 & $[46,46,46,42,42,41,41]$ & 1181 & 2.8054 & 4.789\\
 7 & $[44,44,44,44,44,44,40]$ & 1085 & 2.8066 & 4.787\\
 8 & $[45,45,45,43,42,42,42]$ & 1064 & 2.8066 & 4.788\\
 9 & $[45,45,44,43,43,43,41]$ &  974 & 2.8067 & 4.787\\
10 & $[44,44,44,44,44,42,42]$ &  931 & 2.8070 & 4.787\\
11 & $[47,47,47,41,41,41,40]$ &  902 & 2.8037 & 4.788\\
12 & $[45,45,43,43,43,43,42]$ &  439 & 2.8069 & 4.788\\
13 & $[45,44,44,43,43,43,42]$ &  307 & 2.8070 & 4.788\\
14 & $[46,45,45,42,42,42,42]$ &  236 & 2.8063 & 4.789\\
15 & $[44,44,44,43,43,43,43]$ &  163 & 2.8073 & 4.787\\
16 & $[47,47,44,43,41,41,41]$ &  153 & 2.8050 & 4.788\\
17 & $[46,46,44,42,42,42,42]$ &  107 & 2.8062 & 4.789\\
18 & $[48,48,48,40,40,40,40]$ &  101 & 2.8014 & 4.788\\
19 & $[46,46,43,43,42,42,42]$ &   44 & 2.8063 & 4.788\\
20 & $[46,46,45,42,42,42,41]$ &   39 & 2.8058 & 4.788\\
\midrule
   & \textbf{Total}           & \textbf{19866}\\
\bottomrule
\end{tabular}
\end{table}

Observations: (1) $\lambda_1$ lies in $[4.787,4.789]$ across all 20 types.
(2) $H$ ranges from $2.801$ to $2.807$, all near $\log_27\approx2.807$.
(3) All 101 Type-18 solutions have a genuine $\mathrm{Aut}(Q_7)$-automorphism
(coordinate permutation combined with a bitwise XOR, i.e.\ the full
affine group $\mathrm{Aut}(Q_7)\cong(\mathbb{Z}_2)^7\rtimes S_7$, not
merely a coordinate permutation) that nontrivially permutes at least
two of the three tied $48$-edge directions; of these, $46/101$ have a
genuine order-$3$ automorphism cyclically permuting all three
directions. (An earlier check that considered only pure coordinate
permutations, without combining them with the XOR/translation part of
$\mathrm{Aut}(Q_7)$, incorrectly found no such automorphisms in a
20-solution sample; that check was in error and has been redone
exhaustively over all 101 solutions.)

\subsection{Shared structural properties}

\begin{proposition}\label{prop:rigid}
Every solution in the $19\,866$ satisfies:
\begin{enumerate}[label=(\roman*)]
  \item Degree sequence $\{4^{32},5^{96}\}$.
  \item $\lambda_1$ ranges over $[4.78543,4.79129]$ across all $19\,866$
    solutions (exhaustively computed; not a single shared three-decimal
    value), with $\lambda_1+\lambda_n=0$ throughout, automatic for any
    subgraph of the bipartite $Q_7$ (sanity check, not an independent
    finding). The type-wise means of Table~\ref{tab:types} round to
    $4.787$--$4.789$.
  \item $\mathrm{Tr}(A^4)=5216$ ($C_4$-free trace equality).
  \item Local maximality: no edge of $Q_7$ can be added without creating a $C_4$.
  \item Every edge of $Q_7$ appears in at least one solution
    (verified by taking the union of the edge sets over all $19\,866$ solutions).
\end{enumerate}
\end{proposition}

\subsection{Solution landscape}\label{sec:landscape}

Pairwise Hamming distances $d_H(G,G')=|E(G)\triangle E(G')|$, computed
exhaustively over all $\binom{19\,866}{2}\approx1.97\times10^8$ pairs,
range from $\mathbf{6}$ (attained by a pair of profile-type-1 and
profile-type-2 solutions, differing by just 3 edges each way) to
$\mathbf{274}$. A random sample of $5\,000$ pairs gives a mean
$\approx195.4$ and median $196$, both stable across resampling; the
sample-based min and max are not, and should not be read as close to
the true extremes above.

Over $1\,076$ independent trials targeting $305$ edges,
the minimum $C_4$ violation count was $2$ (never $0$).

\begin{conjecture}\label{conj:q7}
$\ex(Q_7,C_4)=304$.
\end{conjecture}

\subsection{Odd-square audit}\label{sec:oddsquareaudit}

We audited all $19\,866$ solutions of Section~\ref{sec:q7} against the
odd-square condition of Section~\ref{sec:oddsquare} (every square meets the
solution in exactly one or three edges). Exactly $389$ solutions are
odd-square; the remaining $19\,477$ are not. Table~\ref{tab:oddsquareaudit}
gives the dimension-profile breakdown of the 389 odd-square solutions; each
belongs to one of only 4 of the 20 profile types of
Table~\ref{tab:types} (ranks 5, 7, 10, and 18 there).

\begin{table}[ht]
\centering
\caption{Dimension-profile breakdown of the 389 odd-square solutions among
  the $19\,866$ published 304-edge $Q_7$ solutions.}
\label{tab:oddsquareaudit}\medskip
\begin{tabular}{lc}
\toprule
Dimension profile (sorted) & Count\\
\midrule
$\{44,44,44,44,44,44,40\}$ & 127\\
$\{44,44,44,44,44,42,42\}$ & 83\\
$\{48,48,48,40,40,40,40\}$ & 101\\
$\{46,46,46,42,42,42,40\}$ & 78\\
\midrule
Total odd-square & 389\\
Non-odd-square (remaining 16 profile types) & 19\,477\\
\bottomrule
\end{tabular}
\end{table}

Since the odd-square subclass accounts for only $389/19\,866\approx2.0\%$ of
the published solutions, the shared structural core of
Proposition~\ref{prop:rigid} and the 20-type classification of
Section~\ref{sec:profiles} are not fully explained or subsumed by
the odd-square construction: the great majority of 304-edge $C_4$-free
subgraphs of $Q_7$ found here lie outside the fully frustrated switching
class of~\cite{DPTV79}. The one exception is Type~18
($[48,48,48,40,40,40,40]$, 101 solutions), all of which are odd-square;
the other three types containing odd-square solutions (ranks 5, 7, and
10) are mixed, containing both odd-square and non-odd-square members.

%% ============================================================
\section{Computational Method}\label{sec:method}
%% ============================================================

The successful outcomes described in this section (the released edge
sets themselves) are machine-verifiable certificates: any reader can
confirm $C_4$-freeness and the reported edge counts directly from the
supplied data with \texttt{verify.py}, independently of how the search
was run. The quantitative claims about the \emph{search process} that
produced them — ``$19\,866$ independent runs'', ``$1\,076$ independent
trials'' at $305$ and at $681$ edges, and the resampling statistics of
Section~\ref{sec:landscape} — describe unarchived search history. The
search program itself, \texttt{c4free\_sa.py}, is released alongside
this manuscript, but the specific run seeds, logs, and driver
invocations of the actual production runs that produced the released
$19\,866$ solutions and the $1\,076$-trial statistics are not supplied.
These process-level claims should be read as the authors' record of the
search that was run, not as independently reproducible certificates in
the same sense as the edge sets; a reader wishing to verify the exact
process-level statistics would need to rerun an equivalent search
(possible with the released code, but not reproducible bit-for-bit
without the original seeds), not merely check the supplied files.

\subsection{Two-phase simulated annealing}

\noindent\textbf{Phase 1 (Penalty SA).}
We minimise $f(E)=-|E|+\lambda V$ where $V$ is the number of $C_4$ violations
and $\lambda\in[0.30,0.90]$.
Temperature schedule: geometric from $T_0\in[0.20,4.00]$ to
$T_1\in[0.001,0.030]$ over $3$--$15$ million steps.
A move toggles a single edge; the change in $f$ is computed incrementally
in $O(\deg)$ time using precomputed $C_4$-to-edge incidence lists.

\noindent\textbf{Phase 2 (Swap SA).}
Edge count is fixed; swap moves (remove one edge, add one non-edge)
minimise $V$. Phase 2 only runs when phase 1 ends with $V\le4$
violations (Section~\ref{sec:method} intro). In the released code, the
step budget for phase 2 is chosen by the phase-1 violation count $v$,
not by dimension: $2\times10^8$ steps if $v\le1$, otherwise
$5\times10^7$ steps. In practice this correlates with dimension, since
phase 1 typically ends closer to $v\le1$ at the $Q_7$/305-edge target
than at the $Q_8$/681-edge target, but the branch itself is on $v$, not
on $n$ or the target edge count.

\noindent\textbf{Diversification.} The released search code
(\texttt{c4free\_sa.py}) computes, before each trial, a random element of
$\mathrm{Aut}(Q_n)$ (random bit permutation composed with random bit
flip) applied to the current best solution, intended as a diversified
restart point. However, this computed value is not passed into
\texttt{phase1\_sa}, which unconditionally initialises from a fresh
uniform-random edge set of the target size
(\texttt{rng.sample(range(ne), target)}) regardless of any prior best
solution; the automorphism-based restart is dead code in the released
implementation. (The script's own docstring originally described this
mechanism without noting the bug; we have appended a note to the
docstring documenting it, without altering any functional code, so
that a reader of the code alone is not misled.) The diversity actually
observed across trials and across the $19\,866$ collected solutions
therefore comes from this per-trial fresh random initialisation
together with per-trial randomised annealing parameters ($\lambda$,
$T_0$, $T_1$, step count, and violation cap, each redrawn independently
every trial), not from
any automorphism-based mechanism. We note for completeness that even
had the intended mechanism been wired correctly, applying an
automorphism to a solution does not change which $\mathrm{Aut}(Q_n)$-orbit
it belongs to, and the annealing moves used here (uniform edge
toggle/swap proposals against an automorphism-invariant objective) admit
an $\mathrm{Aut}(Q_n)$-equivariant Markov kernel; any diversification
benefit would have come only from decorrelating the specific labelled
trajectory taken by a finite stochastic run; it does not diversify the
orbit sampled from.

\subsection{Comparison with known lower bounds}

Brass--Harborth--Nienborg's weaker bound is stated for $n\ge9$; the
values below evaluate that formula at $n=7,8$ as an extrapolation for
illustrative comparison, not as a proven bound at those dimensions
(see Table~\ref{tab:values} and the Remark in
Section~\ref{sec:q8}).
\begin{align*}
n=7:&\quad\tfrac12(7+0.9\sqrt7)\cdot64\approx300.2\text{ (extrapolated)},\quad\text{bound: }304\;(\Delta=+3.8,\,+1.27\%),\\
n=8:&\quad\tfrac12(8+0.9\sqrt8)\cdot128\approx674.9\text{ (extrapolated)},\quad\text{bound: }682\;(\Delta=+7.1,\,+1.05\%).
\end{align*}
The $n=8$ figure is the odd-square reconstruction of
Section~\ref{sec:oddsquare}; the simulated-annealing method of this section
independently reaches 680 edges on $Q_8$ ($\Delta=+5.1$, $+0.75\%$;
Section~\ref{sec:q8}).

\subsection{Regularity trend across dimensions}\label{sec:regularity}

Degree sequences across dimensions: $Q_6$: $\{4^{56},5^{8}\}$;
$Q_7$: $\{4^{32},5^{96}\}$; $Q_8$: $\{4^8,5^{160},6^{88}\}$ (Solution~A;
Section~\ref{sec:q8}). The minimum degree is $4$ in all three cases; the
degree support is $\{4,5\}$ for both $Q_6$ and $Q_7$ (unchanged from
$n=6$ to $n=7$) and widens to $\{4,5,6\}$ only at $Q_8$. Beyond this, the
exponents do not follow an evident closed-form pattern: $Q_6$'s
$\{56,8\}$ split is markedly more skewed than $Q_7$'s $\{32,96\}$ or
$Q_8$'s $\{8,160,88\}$, so we do not claim a general regularity trend
here.

%% ============================================================
\section{Computational Reproduction of $\ex(Q_6,C_4)=132$}\label{sec:q6}
%% ============================================================

Harborth and Nienborg~\cite{HN94} proved $\ex(Q_6,C_4)=132$ by combinatorial
arguments. We reproduce the lower-bound (construction) side fully
independently and computationally (Section~\ref{sec:q6lower}); for the
upper-bound side we give an ILP formulation whose feasible value matches
132 (Section~\ref{sec:q6upper}), but we did not independently verify
its optimality within a practical runtime, so the upper bound
$\ex(Q_6,C_4)\le132$ itself rests on~\cite{HN94}, not on our own solve.

\subsection{Lower bound: explicit construction}\label{sec:q6lower}

Simulated annealing (Section~\ref{sec:method}) finds a $C_4$-free subgraph of
$Q_6$ on 132 edges, certified by exhaustive enumeration of all 240 four-cycles.
The explicit edge list is provided in \texttt{q6\_edges\_132.jsonl}.

\subsection{Upper bound: integer linear program}\label{sec:q6upper}

Label the 192 edges of $Q_6$ as $x_1,\ldots,x_{192}\in\{0,1\}$,
and let $\mathcal{C}_6$ denote the set of all 240 four-cycles of $Q_6$.
For each four-cycle $\{e_1,e_2,e_3,e_4\}\in\mathcal{C}_6$,
the constraint $x_{e_1}+x_{e_2}+x_{e_3}+x_{e_4}\le3$ ensures $C_4$-freeness.
The ILP is:
\[
  \max\sum_{i=1}^{192}x_i \quad\text{subject to}\quad
  \sum_{e\in C}x_e\le3\;\;\forall C\in\mathcal{C}_6,\quad x_i\in\{0,1\}.
\]
This formulation has 192 binary variables and 240 constraints.
The MPS file \texttt{q6\_ilp.mps} is provided for independent verification.

\begin{proposition}\label{prop:q6ilp}
The ILP above has feasible value $132$, matching the released 132-edge
construction (Section~\ref{sec:q6lower}); combined with
Harborth and Nienborg's independent combinatorial proof
that $\ex(Q_6,C_4)\le132$~\cite{HN94}, this gives $\ex(Q_6,C_4)=132$.
We did not independently verify optimality of $132$ for this ILP
instance within a practical runtime (see below); the upper bound
$\ex(Q_6,C_4)\le132$ used here rests on~\cite{HN94}, not on our own
solve of the ILP.
\end{proposition}

We independently re-solved \texttt{q6\_ilp.mps} with HiGHS~1.15.1
(Python bindings via \texttt{highspy}; default settings, single thread,
single-threaded x86\_64 Linux VM, Intel Xeon @ 2.10GHz, 4GB RAM;
invocation: \texttt{Highs().readModel(...)} then \texttt{run()} with
\texttt{time\_limit} set to 60 and 260 seconds respectively). A
feasible solution matching the released 132-edge construction was
found quickly, but the solver did not close the optimality gap within
several minutes: after 60s the gap was 6.82\% (dual bound $-141$), and
after 260s it was 6.06\% (dual bound $-140$), against the primal bound
of $-132$ throughout. As these figures are implementation- and
machine-dependent, they should be read as evidence that the gap does
not close quickly under a generic default-settings solve, not as a
precise benchmark. This is
consistent with the large automorphism group of $Q_6$
($|\mathrm{Aut}(Q_6)|=6!\cdot2^6=46\,080$) producing many symmetric
alternative optima, which is known to slow generic branch-and-bound
without dedicated symmetry-breaking. We do not claim generic MIP
solvers verify this instance's optimality quickly; the 132-edge value
should be regarded as independently established by
Harborth--Nienborg's combinatorial proof~\cite{HN94}, with the ILP
providing a machine-checkable feasible witness for the primal side and
a formulation that a sufficiently long or specialised solve can, in
principle, verify on the dual side.

%% ============================================================
\section{Open Problems}\label{sec:open}
%% ============================================================

\begin{enumerate}
  \item Prove or disprove $\ex(Q_7,C_4)=304$ (Conjecture~\ref{conj:q7}).
  \item\label{op:q8exact} Determine $\ex(Q_8,C_4)$ exactly. Is it $682$, or can a
    non-odd-square construction exceed $g(8)=682$? The $Q_7$ audit of
    Section~\ref{sec:oddsquareaudit} shows the two questions are genuinely
    different for $n=7$ (the general extremal value $304$ is attained
    overwhelmingly by non-odd-square solutions), so no assumption is made
    here about which way $n=8$ resolves.
  \item What is the number of $\mathrm{Aut}(Q_7)$-orbits among all 304-edge
    $C_4$-free subgraphs of $Q_7$? Among the 389 odd-square solutions
    specifically?
  \item Is the degree sequence $\{4^{32},5^{96}\}$ necessary for any locally
    maximal $C_4$-free subgraph of $Q_7$ on 304 edges?
  \item Can the flag algebra method of~\cite{Bab12,BHLL12} yield tight bounds
    for small $n\le8$?
  \item Does the degree support of locally maximal $C_4$-free subgraphs
    of $Q_n$ near the extremal edge count widen beyond $\{4,5\}$ again
    for $n\ge9$ (as it did once, from $Q_7$ to $Q_8$; Section~\ref{sec:regularity}),
    and is minimum degree $4$ necessary at these edge counts for
    $n\ge9$?
  \item We now know $g(6)=\ex(Q_6,C_4)=132$ (Theorem~\ref{thm:oddsquare}
    combined with~\cite{HN94}). For $n=7,8$, Theorem~\ref{thm:oddsquare}
    only establishes $g(7)=304$ and $g(8)=682$; whether
    $\ex(Q_7,C_4)=g(7)$ and $\ex(Q_8,C_4)=g(8)$ remains exactly as open
    as Conjecture~\ref{conj:q7} and Open Problem~\ref{op:q8exact} above, since
    $\ex(Q_n,C_4)$ itself is not known exactly for $n=7,8$. (The
    coincidence at $n=6$ is in any case not explained by the
    odd-square construction alone for $n=7$: only 389 of the
    $19\,866$ known $304$-edge $Q_7$ solutions are odd-square.) Does
    $g(n)=\ex(Q_n,C_4)$ continue to hold for $n=9,10,\ldots$ (once
    $\ex(Q_n,C_4)$ is itself determined), or does the gap
    $\ex(Q_n,C_4)-g(n)$ become positive at some $n$?
\end{enumerate}

\section*{Acknowledgements}
The author thanks B.~Lidick\'y (Iowa State University) for the endorsement
and for pointing out the reference~\cite{Bab12}.
The author thanks S.~Lai for identifying the relevance
of~\cite{DPTV79,MPR95} to the $Q_7$ and $Q_8$ constructions of this paper
and for independently checking the reconstructed odd-square material of
Section~\ref{sec:oddsquare}.
The author also thanks the open-source community for the HiGHS and SCIP
solvers used in ILP experiments.

\noindent\textit{Use of generative AI.} The author used Claude (Anthropic)
for: (i)~language editing of the abstract; (ii)~assistance in writing the
supplementary verification script (\texttt{verify.py}) used to confirm
$C_4$-freeness of the released edge-list certificates; and, for this
revision, (iii)~retrieval and cross-checking of the primary sources cited
in Section~\ref{sec:oddsquare} and (iv)~drafting assistance for the revised
text. All definitions, constructions, theorems, computations, and analyses
are the author's own, including the 256-bit spin configuration used in
the $Q_8$ witness of Section~\ref{sec:oddsquareverify}, whose derivation
is documented there. The author reviewed all
AI-assisted output, independently executed the verification and witness
scripts against the released data, and takes full responsibility for the
accuracy and integrity of the work. The specific model version(s) of
Claude used across the original submission and this revision were not
recorded at the time and cannot be confirmed retrospectively; the
disclosure above is accurate as to which tasks AI assistance was used
for, but not as to the exact model version for each.

\section*{Disclosure statement}
No potential conflict of interest was reported by the author.

\begin{thebibliography}{99}
\bibitem{Bab12}
R.~Baber,
Tur\'an densities of hypercubes,
arXiv:1201.3587, 2012.

\bibitem{BHLL12} J.~Balogh, P.~Hu, B.~Lidick\'y, and H.~Liu, Upper bounds on
  the size of 4- and 6-cycle-free subgraphs of the hypercube,
  \textit{Eur.\ J.\ Combin.}\ \textbf{35} (2014), 75--85.
\bibitem{BHN95} P.~Brass, H.~Harborth, and H.~Nienborg, On the maximum number
  of edges in a $C_4$-free subgraph of $Q_n$,
  \textit{J.\ Graph Theory}\ \textbf{19} (1995), 17--23.
\bibitem{Chu92} F.~Chung, Subgraphs of a hypercube containing no small even
  cycles, \textit{J.\ Graph Theory}\ \textbf{16} (1992), 273--286.
\bibitem{DPTV79} B.~Derrida, Y.~Pomeau, G.~Toulouse, and J.~Vannimenus,
  Fully frustrated simple cubic lattices and the overblocking effect,
  \textit{J.\ Physique}\ \textbf{40} (1979), 617--626.
\bibitem{EMR92} M.\,R.~Emamy-K., P.~Guan, and P.\,I.~Rivera-Vega, On the
  characterization of the maximum squareless subgraphs of the 5-cube,
  \textit{Congr.\ Numer.}\ \textbf{88} (1992), 97--109.
\bibitem{Erd84} P.~Erd\H{o}s, On some problems in graph theory, combinatorial
  analysis and combinatorial number theory, in: \textit{Graph Theory and
  Combinatorics} (B.~Bollob\'as, ed.), Academic Press, 1984, pp.~1--17.
\bibitem{HN94} H.~Harborth and H.~Nienborg, Maximum number of edges in a
  six-cube without four-cycles, \textit{Bull.\ ICA}\ \textbf{12} (1994), 55--60.
\bibitem{LNSW25} S.~Laplante, R.~Naserasr, A.~Sunny, and Z.~Wang,
  Sensitivity conjecture and signed hypercubes,
  \textit{ACM Trans.\ Comput.\ Theory}, 2025,
  \url{https://doi.org/10.1145/3777401}
  (preprint: HAL hal-04080411).
\bibitem{MPR95} E.~Marinari, G.~Parisi, and F.~Ritort, The fully frustrated
  hypercubic model is glassy and aging at large $D$,
  \textit{J.\ Phys.\ A}\ \textbf{28} (1995), 327--334, arXiv:cond-mat/9410089.
\bibitem{TW09} A.~Thomason and P.~Wagner, Bounding the size of square-free
  subgraphs of the hypercube, \textit{Discrete Math.}\ \textbf{309} (2009),
  1730--1735.
\end{thebibliography}
\end{document}
